\documentclass[12pt]{article}
\usepackage[utf8]{inputenc}
\usepackage[T1]{fontenc}
\usepackage[margin=1in]{geometry}
\usepackage{setspace}
\usepackage{xcolor}
\usepackage[colorlinks=true, urlcolor=blue!60!black, linkcolor=blue!60!black]{hyperref}
\setlength{\parindent}{0em}
\setlength{\parskip}{0.6em}
\renewcommand{\baselinestretch}{1.0}
\begin{document}
\begin{center}
{\LARGE \textit{Ashcroft v.\ Free Speech Coalition}, 535 U.S.\ 234 (2002)\par}
\vspace{0.5em}
{\large Case Brief --- Two-Page Edition\par}
\end{center}

\section{Memory Jogger}
The Court struck down the CPPA's bans on images that ``appear to be'' of minors or ``convey the impression'' of minors as substantially overbroad: government may not suppress non-obscene speech made without harming real children (\href{https://supreme.justia.com/cases/federal/us/535/234/}{535 U.S.\ 234}).

\section{Detailed Case Facts}
The Child Pornography Prevention Act of 1996 extended the federal ban beyond images of actual minors to any image that ``appears to be'' of a minor in sexually explicit conduct (\S~2256(8)(B)), reaching computer-generated images and youthful-looking adults, and to any depiction marketed so that it ``conveys the impression'' of a minor (\S~2256(8)(D)). Neither required obscenity. Congress cited grooming, appetite-whetting, market effects, and proof problems. An adult-entertainment trade group and several artists brought a pre-enforcement facial challenge.

\section{Procedural History}
The district court upheld the Act; the Ninth Circuit reversed. Four other circuits had upheld it, and the Supreme Court granted certiorari.

\section{Judicial Votes}
Affirmed, 6--3 on virtual images and 7--2 on youthful adults and pandering. Kennedy wrote for Stevens, Souter, Ginsburg, and Breyer; Thomas concurred in the judgment; O'Connor concurred in part and dissented in part (Rehnquist and Scalia joining Part II); Rehnquist dissented, joined mostly by Scalia.

\section{Holding}
Both provisions are facially overbroad. Speech that is neither obscene under \textit{Miller} nor produced by exploiting real children under \textit{Ferber} keeps First Amendment protection, and Congress may not ban it because it might indirectly encourage abuse or complicate prosecutions.

\section{Analysis of Opinions}
Kennedy held that the Act ignored \textit{Miller}, threatening works like \textit{Romeo and Juliet} and \textit{American Beauty}, and that \textit{Ferber} rested on speech ``intrinsically related'' to real abuse. Indirect-harm theories failed under \textit{Stanley} and \textit{Brandenburg}; banning protected speech because it resembles unprotected speech ``turns the First Amendment upside down''; and the pandering clause tainted innocent downstream possessors. Thomas left room for a narrow law if technology made real cases unprosecutable. O'Connor would uphold the ban on images ``virtually indistinguishable from'' real children; Rehnquist would have saved the whole Act by narrowing construction.

\section{Examples: Future Applications}
\textbf{Same-Side:} a state ban on non-obscene animated depictions of fictional teenagers falls. \textbf{Opposite-Side:} a prosecution for obscene virtual images under \S~1466A, which incorporates \textit{Miller}, is valid. \textbf{Fence-Sitter:} photorealistic generative-AI images, especially from models trained on real abuse material, test whether output is ``intrinsically related'' to real abuse.

\section{Critique}
Child-protection scholars argue the production-harm test understates real dangers, a critique sharpened by generative AI. Doctrinal critics fault the Court for using overbreadth rather than avoidance or severability, while Joseph Blocher and others see the case as reinforcing First Amendment categoricalism.

\section{Key Quotations}
\begin{quote}
``In contrast to the speech in \textit{Ferber}, speech that itself is the record of sexual abuse, the CPPA prohibits speech that records no crime and creates no victims by its production. Virtual child pornography is not `intrinsically related' to the sexual abuse of children, as were the materials in \textit{Ferber}.''
\end{quote}
\noindent --- Justice Kennedy, majority opinion, 535 U.S.\ at 250.

\end{document}
